%% notes-tex/025-s3-closure/sections/3-strength.tex

\section{Strength: the score against its own definition\secstatus{todo}}
\label{sec:strength}

\tcfigfit[0.98]{figures/s3_closure_strength.pdf}{Recomputed against stored
interaction scores. (a) Within one source screen: each Costanzo 2016 row's
$f_{ab} - f_a f_b$ against the score reported on the same row, over the
1{,}227{,}294 rows whose gene pair is a closure pair. (b) In the build:
Eq.~\ref{eq:eps} from the build's single and double fitness against the stored
digenic score, all 728{,}856 closure doubles with both. (c) In the build:
Eq.~\ref{eq:tau} from the build's fitness against the stored trigenic score, the
352{,}505 triples with full closure. Hexagon shade is log count; the dashed line is
identity.}{fig:strength}

\begin{table}[htbp]
\centering
\footnotesize
\caption{Recomputed against stored score, by stratum. Within-screen rows recompute
from the columns of one raw source row; every other row recomputes from the build's
stored fitness. ``Measured mean restored'' takes the essentiality 0 back out of each
tainted single by $k/(k-1)$ (Section~\ref{sec:hazards}). The best $r$ is bold.
Generated by \texttt{s3\_closure\_recompute.py}.}
\label{tab:strength}
%% SOURCE: experiments/025-solid-growth/scripts/s3_closure_recompute.py (results/s3_closure_recompute_summary.json) -- GENERATED, do not edit

\begin{tabular}{lrrrrr}
\toprule
stratum & $n$ & $r$ & $\rho$ & slope & rmse \\
\midrule
within one screen, Costanzo 2016 (raw row) & 1,227,294 & 0.999 & 1.000 & 1.00 & 0.002 \\
within one screen, Kuzmin 2018 (raw row) & 197,876 & \textbf{1.000} & 1.000 & 1.00 & 0.000 \\
within one screen, Kuzmin 2020 (raw row) & 601,330 & 0.996 & 0.997 & 1.00 & 0.004 \\
digenic, all closure doubles & 728,856 & 0.445 & 0.493 & 0.99 & 0.072 \\
digenic, one screen only & 36,499 & 0.339 & 0.314 & 0.95 & 0.120 \\
digenic, merged screens & 692,357 & 0.459 & 0.509 & 0.99 & 0.068 \\
digenic, no essential single & 589,395 & 0.598 & 0.633 & 0.98 & 0.042 \\
digenic, an essential single (stored mean) & 139,461 & 0.452 & 0.382 & 0.97 & 0.140 \\
digenic, an essential single (measured mean restored) & 139,461 & 0.491 & 0.441 & 0.96 & 0.071 \\
trigenic, all triples & 352,505 & 0.230 & 0.224 & 0.63 & 0.171 \\
trigenic, one screen only & 339,616 & 0.239 & 0.233 & 0.66 & 0.169 \\
trigenic, merged screens & 12,889 & 0.114 & 0.153 & 0.30 & 0.223 \\
trigenic, no essential single & 278,060 & 0.205 & 0.211 & 0.58 & 0.169 \\
trigenic, an essential single (stored mean) & 74,445 & 0.255 & 0.188 & 0.60 & 0.180 \\
trigenic, an essential single (measured mean restored) & 74,445 & 0.250 & 0.186 & 0.60 & 0.167 \\
trigenic, from stored dmi & 350,318 & 0.213 & 0.174 & 0.48 & 0.144 \\
\bottomrule
\end{tabular}

\end{table}

\notesec{The formula is exact inside a screen}
Figure~\ref{fig:strength}a and the first three rows of Table~\ref{tab:strength} are
the control. On the raw Costanzo 2016 rows the identity reproduces the reported score
at $r = 0.999$, slope 1.00, root-mean-square error 0.002; on the raw Kuzmin 2018 rows
at $r = 1.000$ with zero residual; on the raw Kuzmin 2020 rows at $r = 0.996$. The
Kuzmin 2020 figure needs one substitution to reach that: 99.6 percent of its digenic
rows carry no query single-mutant fitness, and the identity reproduces the reported
score only when that missing value is set to 1.0, which is what the supplementary
text says the scoring did (``all NaN fitness estimates were assigned a value of
1.0'', \texttt{si/si1.md} line 41). Hazard H5 is therefore not a hypothesis: the
reported Kuzmin 2020 digenic scores were computed against a query fitness of 1.0, and
the build inherits scores that no stored fitness can reproduce for those pairs. The
small vertical streak at a reported score of 0 in Figure~\ref{fig:strength}a is the
2.1 percent of Costanzo rows with a missing query fitness, where the same rule
applies.

\notesec{Digenic, in the build}
The same identity on the build's own values reaches $r = 0.445$, $\rho = 0.493$,
slope 0.99 and a root-mean-square error of 0.072 over 728{,}856 doubles
(Figure~\ref{fig:strength}b). The slope says the identity holds in expectation; the
error says it fails in value, by an amount larger than the stored score's own spread.
The strata separate the causes. Doubles whose two singles carry no essentiality entry
reach $r = 0.598$ with error 0.042; doubles touching an essentiality-tainted single
reach $r = 0.452$ with error 0.140 and an intercept of 0.116, which is the signature
of a product $f_a f_b$ that is too small because one factor was averaged with a 0.
Restoring the measured mean of the tainted single moves that stratum to $r = 0.491$,
intercept 0.014, error 0.071: half the error of that stratum was the essentiality 0.
Doubles kept from one screen reach only $r = 0.339$, against $r = 0.459$ for merged
doubles, because a merged double's fitness and score are both means over screens and
average some of the same noise away, while a single-screen double's fitness is
compared against singles averaged over other screens.

\notesec{Trigenic, in the build}
Eq.~\ref{eq:tau} on the build's fitness reaches $r = 0.230$, slope 0.63 and error
0.171 over the 352{,}505 triples whose three singles and three doubles are all present
(Figure~\ref{fig:strength}c). The error is 2.7 times the standard deviation of the
stored trigenic score, so the reconstruction carries more propagated measurement
error than signal, which is the attenuation the 2026-09-04 recapitulation computed.
The essentiality strata behave as in the digenic case (intercept $-0.085$ with the
stored singles, $0.022$ with the measured mean restored) without moving $r$, because
five terms of noise remain. Triples merged from two queries reach $r = 0.114$, and
using the stored digenic scores in place of the recomputed products gives
$r = 0.213$, no rescue. At the intermediate threshold of $|\tau| > 0.08$ the build's
fitness calls 169{,}476 triples where the stored scores call 46{,}889, and the two
agree on 29{,}327, a recall of 0.63 and a precision of 0.17.

\notesec{The trigenic formula is exact inside a screen, and one term carries it}
The digenic control above fits on one raw row. The trigenic identity does not, because
the supplementary methods place its terms in three different measurements
(\texttt{kuzminSystematicAnalysisComplex2018/si/si1.md} line 191): the double-mutant
query fitness $f_{ij}$, the array single fitness $f_k$ and the triple fitness
$f_{ijk}$ sit on the trigenic row itself, while $\varepsilon_{ik}$, $\varepsilon_{jk}$
and the query singles $f_i$, $f_j$ come from the two single-mutant control query
strains, gene plus the \emph{HO} deletion, screened against the same array strain in
the same batch. Matching those rows by array strain identifier recovers the released
score exactly, and the form that does it names what the source used:
\[
  \tau_{ijk} = f_{ijk} - f_{ij} f_k - \varepsilon_{ik} - \varepsilon_{jk},
\]
that is, Eq.~\ref{eq:tau} with the single-mutant query fitness set to 1. This reproduces
the published adjusted score on 95.0 percent of the 91{,}111 Kuzmin 2018 trigenic rows
and 92.0 percent of the 301{,}798 Kuzmin 2020 rows to within $10^{-4}$, the rounding of
the published five-decimal value, at a median absolute residual of $2.9\times10^{-5}$
(Table~\ref{tab:trigenic-within}). The model in Eq.~\ref{eq:tau} is therefore the right
one and is implemented correctly here; nothing in the published correction is missing.

The substitution of 1 for the query singles is a property of the source, not an
approximation made here. Kuzmin 2020 releases a query fitness on 0.4 percent of its
control rows, and its supplementary methods state that every missing fitness was assigned
1.0 during scoring. Kuzmin 2018 releases one on 99.2 percent of its control rows and did
not use them: reading those measured values into $\varepsilon_{ik} f_j$ and
$\varepsilon_{jk} f_i$ raises the median residual from $2.9\times10^{-5}$ to
$1.6\times10^{-3}$, fifty times worse. What a control screen contributes to a trigenic
score is therefore its interaction, not its fitness.

\begin{table}[htbp]
\centering
\footnotesize
\caption{The trigenic identity inside one screen, matched by array strain identifier.
The first row of each pair takes every term from the screen that measured it. The
second drops the double-mutant query fitness $f_{ij}$, which is the only change made,
and is what a closure over the database is forced into because that measurement is not
a record. Generated by \texttt{s3\_closure\_trigenic\_within\_screen.py}.}
\label{tab:trigenic-within}
%% SOURCE: experiments/025-solid-growth/scripts/s3_closure_trigenic_within_screen.py (results/s3_closure_trigenic_within_screen.json) -- GENERATED, do not edit

\begin{tabular}{llrrrrr}
\toprule
screen & terms & $n$ & $r$ & $\rho$ & slope & rmse \\
\midrule
Kuzmin 2018 & every term from the same screen & 91,111 & \textbf{0.985} & 0.971 & 0.98 & 0.009 \\
\quad & the query double fitness replaced by $f_i f_j$ & 91,111 & 0.538 & 0.433 & 1.03 & 0.088 \\
Kuzmin 2020 & every term from the same screen & 301,706 & 0.976 & 0.975 & 0.97 & 0.014 \\
\quad & the query double fitness replaced by $f_i f_j$ & 301,706 & 0.419 & 0.490 & 0.97 & 0.151 \\
\bottomrule
\end{tabular}

\end{table}

That leaves one term to explain the build. Dropping $f_{ij}$ alone, and changing nothing
else, collapses the same reconstruction to $r = 0.495$ on Kuzmin 2018 and $r = 0.420$ on
Kuzmin 2020. The double-mutant query fitness is the
term that carries the trigenic score, which is the same statement as saying that the
digenic interaction between the two query genes is large and cannot be inferred from
their single-mutant fitness. The build has no record of it: across the whole Kuzmin
2018 screen only 172 query pairs carry a released double-mutant query fitness, 129 of
them appear as a double in the 029 build, and those records are a different
measurement of the same genotype, an array screen by Costanzo or by Kuzmin's own
digenic panel, agreeing with the query-strain value at $r = 0.777$ with a median
absolute difference of 0.045. That difference enters $\tau$ multiplied by $f_k \approx 1$,
so it is more than half the 0.08 calling threshold on its own.
